%
% mail-refile-cancel-dup.tex  (FROZEN NEGATIVE CONTROL: do not edit)
%
% ReadCancelled leaves the filed messages in the outbox (the release after
% the loop) and the fetch does not reconcile a filed copy with its original.
% Expected: the goal cancelled=ztrue & refileTimedOut=zfalse & a message in
% both stores is reached right after ReadCancelled; MailInv is violated at the
% next fetch.
%
% Z Specification for a partial read on a consuming relay: the refile of the
% remainder, and the confirmation of the batch (DES-095, DES-099).
%
% Companion to session-model.tex and nats-relay.tex.  session-model.tex keeps
% the unread mail of a session as a count ($durable$, $shown$); this model
% needs the messages themselves, because the properties are about which
% message is where.  It is kept apart because those two specifications
% already have large state spaces.
%
% Type-check: fuzz -t mail-refile.tex
% Model-check: see the end of this file.
%

\documentclass[a4paper,10pt,fleqn]{article}
\usepackage[margin=1in]{geometry}
\usepackage{fuzz}
\usepackage[colorlinks=true,linkcolor=blue,citecolor=blue,urlcolor=blue]{hyperref}

\begin{document}

\title{Biff Mail Refile and Confirmation: A Z Specification}
\author{Punt Labs}
\date{November 2026}
\maketitle

%%----------------------------------------------------------------------
\section{Introduction}
%%----------------------------------------------------------------------

The NATS relay \emph{consumes} mail when it is fetched: a fetch acknowledges
the messages it returns, and the work queue deletes what is acknowledged.
After a fetch the mail exists only in the memory of whoever fetched it.

The daemon's \texttt{read} fetches a bounded batch, fits as much of it into
the reply as the read budget allows, and is left with a remainder.  The
remainder used to stay in the daemon's outbox, out of reach of everything else
that drains a session: ending the session, or a read from another process,
found the relay empty.  A test on a real relay found 0 of 26 messages
reachable after a partial read.  The remainder is now filed back into the
relay under an id derived from its own (\texttt{Message.refiled}), so that
filing it twice cannot make two copies.

The reply carries the batch.  The batch stays in the outbox until the client
confirms it, and the client confirms only after it has printed it
(\texttt{request\_held}, \texttt{HeldReply.confirm}).  A confirmation sent
before the print let a client that died in between mark as read a message it
had never shown.

This specification states what is claimed and what is not.  The claims hold
when nothing fails; each failure that voids one is an explicit operation.

\begin{enumerate}
  \item \textbf{No remainder is lost.} After any partial read, every message
    the batch does not carry is in the relay or, if the relay refused it,
    in the outbox of a session that has not ended.
  \item \textbf{Nothing is duplicated.} A message is held by at most one of
    the relay and the outbox, and the relay holds at most one copy of it.
  \item \textbf{\texttt{more} counts each remaining message once.}
  \item \textbf{A message is marked read only after it was shown.}
  \item \textbf{Nothing vanishes unaccounted.} Every issued message is shown,
    drained, held, or recorded as lost by a named event.
\end{enumerate}

%%----------------------------------------------------------------------
\section{Basic Types}
%%----------------------------------------------------------------------

\begin{zed}
  ZBOOL ::= ztrue | zfalse
\end{zed}

\begin{zed}
  ReadPhase ::= rdIdle | rdFetched
\end{zed}

\begin{axdef}
  fetchMax : \nat \\
  batchMax : \nat
  \where
  fetchMax = 2 \\
  batchMax = 1
\end{axdef}

$fetchMax$ is the bound of one consuming fetch; $batchMax$ is the number of
messages the read budget fits into one reply.  Both are small so that
ProB can enumerate them, and $batchMax < fetchMax$ is what makes a read
partial.

\begin{zed}
  MAIL == 1 \upto 2 \\
  GEN == 0 \upto 1 \\
  WIRE == MAIL \cross GEN
\end{zed}

$MAIL$ is a message as the sender knows it.  $WIRE$ is a message as the relay
knows it: the message and the generation of its id.  A refile gives the
message the next id in its chain; the id is \texttt{uuid5} of the previous
one, so the chain never meets itself.  The model counts the generations
modulo two, which is sound because a message has one live copy (property 2),
and the one pair that would collide is exactly the case the derivation is
there for: a copy filed twice from the same predecessor.

\begin{axdef}
  refile : WIRE \fun WIRE
  \where
  refile = \{ m : MAIL; g : GEN @ (m, g) \mapsto (m, (g + 1) \mod 2) \}
\end{axdef}

\begin{axdef}
  mailOf : \power WIRE \fun \power MAIL
  \where
  \forall s : \power WIRE @ mailOf~s = \{ w : s @ first~w \}
\end{axdef}

%%----------------------------------------------------------------------
\section{State}
%%----------------------------------------------------------------------

\begin{schema}{Store}
  relay : \power WIRE \\
  held : \power WIRE \\
  sent : \power WIRE \\
  phase : ReadPhase
\end{schema}

$relay$ is the session's inbox in the relay, shared with every process that
drains it.  $held$ is the daemon's outbox: process memory, lost with the
process.  $sent$ is the batch of the reply the client has not yet confirmed;
a message in $sent$ is held.  $phase$ is $rdFetched$ between the fetch of a
read and its reply.

\begin{schema}{Client}
  shown : \power MAIL \\
  done : \power MAIL \\
  dupShown : ZBOOL
\end{schema}

$shown$ is what the client has printed, $done$ what has been marked read, and
$dupShown$ records that a reply offered the client a message it had already
printed.

\begin{schema}{Ledger}
  issued : \power MAIL \\
  drained : \power MAIL \\
  lostBatch : \power MAIL \\
  lostRest : \power MAIL
\end{schema}

A ghost.  $issued$ is every message ever filed; $drained$ is what an ending
session moved into its owner's mailbox, where it is reachable.  $lostBatch$
and $lostRest$ are what a failure took from the outbox without the client
having been shown it: the batch of an unconfirmed reply, and the rest.

\begin{schema}{Flags}
  faulted : ZBOOL \\
  cancelled : ZBOOL \\
  refileTimedOut : ZBOOL \\
  moreOut : 0 \upto 2 \\
  moreTruth : 0 \upto 2 \\
  moreSlack : 0 \upto 2
\end{schema}

$faulted$ is set by every operation that is a failure and not a protocol
step: a daemon crash, a relay that refuses a refile, a session that ends
during a read.  A read cancelled half-way and a refile that landed but timed
out are events of their own, $cancelled$ and $refileTimedOut$, so that a
witness query can isolate each; like $faulted$ they leave unfiled mail in the
outbox only, and the claim that the remainder is reachable is stated while
none of the three has happened.  $moreOut$ is the \texttt{more}
the last reply carried, computed as the code computes it; $moreTruth$ is the
number of messages that reply left, computed from the ledger.

\begin{schema}{MailState}
  Store \\
  Client \\
  Ledger \\
  Flags
\end{schema}

\begin{schema}{MailInv}
  MailState
  \where
  sent \subseteq held \\
  done \subseteq shown \\
  shown \subseteq issued \\
  issued \subseteq shown \cup drained \cup mailOf~relay \cup mailOf~held \cup
    lostBatch \cup lostRest \\
  ( faulted = zfalse \land cancelled = zfalse \land refileTimedOut = zfalse )
    \implies lostRest = \emptyset \\
  ( \forall m : MAIL @ \# \{ w : relay | first~w = m \} \leq 1 \land
    \# \{ w : held | first~w = m \} \leq 1 ) \\
  ( \forall w : relay; h : held | first~w = first~h @ w = refile~h ) \\
  moreTruth \leq moreOut \\
  moreOut \leq moreTruth + moreSlack
\end{schema}

The five properties, in the order of the introduction (the first is the
fifth line, the last the second).

\begin{schema}{Invariant}
  MailInv
\end{schema}

\begin{schema}{Init}
  MailState'
  \where
  relay' = \emptyset \\
  held' = \emptyset \\
  sent' = \emptyset \\
  phase' = rdIdle \\
  shown' = \emptyset \\
  done' = \emptyset \\
  dupShown' = zfalse \\
  issued' = \emptyset \\
  drained' = \emptyset \\
  lostBatch' = \emptyset \\
  lostRest' = \emptyset \\
  faulted' = zfalse \\
  cancelled' = zfalse \\
  refileTimedOut' = zfalse \\
  moreOut' = 0 \\
  moreTruth' = 0 \\
  moreSlack' = 0
\end{schema}

%%----------------------------------------------------------------------
\section{Operations}
%%----------------------------------------------------------------------

\begin{schema}{Arrive}
  \Delta Store \\
  \Xi Client \\
  \Delta Ledger \\
  \Xi Flags \\
  m? : MAIL
  \where
  m? \notin issued \\
  relay' = relay \cup \{ (m?, 0) \} \\
  held' = held \\
  sent' = sent \\
  phase' = phase \\
  issued' = issued \cup \{ m? \} \\
  drained' = drained \\
  lostBatch' = lostBatch \\
  lostRest' = lostRest
\end{schema}

A message is filed in the session's inbox.

\begin{schema}{ReadFetch}
  \Delta Store \\
  \Xi Client \\
  \Xi Ledger \\
  \Xi Flags \\
  fetched? : \power WIRE
  \where
  phase = rdIdle \\
  fetched? \subseteq relay \\
  \# fetched? \leq fetchMax \\
  \# relay \leq fetchMax \implies fetched? = relay \\
  relay' = relay \setminus fetched? \\
  held' = held \cup fetched? \\
  sent' = sent \\
  phase' = rdFetched
\end{schema}

The consuming fetch: what it returns leaves the relay and is held.  A
bounded fetch takes the whole inbox when the inbox fits.  A fetched message
that is the refile of a held one is that message filed again, and replaces
it (\texttt{Outbox.merge}); this is what makes a refile that landed but was
reported failed, or a read interrupted before the release, harmless to the
count of copies.

\begin{schema}{ReadReply}
  \Delta Store \\
  \Delta Client \\
  \Xi Ledger \\
  \Delta Flags \\
  batch? : \power WIRE
  \where
  phase = rdFetched \\
  batch? \subseteq held \\
  shown' = shown \\
  done' = done \\
  dupShown' = ( \IF mailOf~batch? \cap shown \neq \emptyset \THEN ztrue \ELSE dupShown ) \\
  \# batch? \leq batchMax \\
  \# held \leq batchMax \implies batch? = held \\
  \# held > batchMax \implies \# batch? = batchMax \\
  relay' = relay \cup refile \limg held \setminus batch? \rimg \\
  held' = batch? \\
  sent' = batch? \\
  phase' = rdIdle \\
  cancelled' = cancelled \\
  refileTimedOut' = refileTimedOut \\
  faulted' = faulted \\
  moreOut' = \# relay' \\
  moreTruth' = \# ( mailOf~( relay' \cup ( held' \setminus batch? ) ) \setminus
    mailOf~batch? ) \\
  moreSlack' = \# \{ h : held' | refile~h \in relay' \}
\end{schema}

The reply and the refile.  The batch is the part of the outbox the budget
fits.  The remainder is filed back, each message under its derived id, and
released from the outbox; the batch stays, unconfirmed.  The reply says
$moreOut' = \# relay'$: \texttt{left = len(pending) - len(batch) - filed +
unfetched}, where $filed$ is the remainder and $unfetched$ the relay's own
count after the refile, so the remainder is subtracted from the one count and
added by the other.  $moreTruth'$ is the number of \emph{distinct messages} the
reply leaves: those held anywhere for a later read, less the ones this
reply carries.  The two agree when no message has two copies; that is the
sense in which \texttt{more} counts each remaining message once.

\begin{schema}{ReadReplyRefused}
  \Delta Store \\
  \Delta Client \\
  \Xi Ledger \\
  \Delta Flags \\
  batch? : \power WIRE
  \where
  phase = rdFetched \\
  batch? \subseteq held \\
  shown' = shown \\
  done' = done \\
  dupShown' = ( \IF mailOf~batch? \cap shown \neq \emptyset \THEN ztrue \ELSE dupShown ) \\
  \# batch? \leq batchMax \\
  \# held \leq batchMax \implies batch? = held \\
  \# held > batchMax \implies \# batch? = batchMax \\
  relay' = relay \\
  held' = held \\
  sent' = batch? \\
  phase' = rdIdle \\
  cancelled' = cancelled \\
  refileTimedOut' = refileTimedOut \\
  faulted' = ztrue \\
  moreOut' = \# ( held \setminus batch? ) + \# relay \\
  moreTruth' = \# ( mailOf~( relay \cup ( held \setminus batch? ) ) \setminus
    mailOf~batch? ) \\
  moreSlack' = \# \{ h : held | refile~h \in relay \}
\end{schema}

The relay refused the refile (\texttt{RELAY\_ERRORS}, logged).  The remainder
stays in the outbox and is offered again by the next read.  The reply still
counts it: \texttt{left = len(pending) - len(batch) - filed + unfetched} with
$filed = 0$, that is the held remainder plus the relay's own backlog, which is
$moreTruth'$; the count invariant is not switched off by a refusal.  Here
$moreSlack' = \# \{ h : held | refile~h \in relay \}$ counts the held messages
whose filed copy is already in the relay, which a refusal never adds.

\begin{schema}{ReadReplyTimedOut}
  \Delta Store \\
  \Delta Client \\
  \Xi Ledger \\
  \Delta Flags \\
  batch? : \power WIRE
  \where
  phase = rdFetched \\
  batch? \subseteq held \\
  shown' = shown \\
  done' = done \\
  dupShown' = ( \IF mailOf~batch? \cap shown \neq \emptyset \THEN ztrue \ELSE dupShown ) \\
  \# batch? \leq batchMax \\
  \# held \leq batchMax \implies batch? = held \\
  \# held > batchMax \implies \# batch? = batchMax \\
  relay' = relay \cup refile \limg held \setminus batch? \rimg \\
  held' = held \\
  sent' = batch? \\
  phase' = rdIdle \\
  cancelled' = cancelled \\
  refileTimedOut' = ztrue \\
  faulted' = faulted \\
  moreOut' = \# ( held \setminus batch? ) + \# relay' \\
  moreTruth' = \# ( mailOf~( relay' \cup ( held \setminus batch? ) ) \setminus
    mailOf~batch? ) \\
  moreSlack' = \# \{ h : held | refile~h \in relay' \}
\end{schema}

The refile \emph{landed} and was reported as a timeout (\texttt{TimeoutError}
is in \texttt{RELAY\_ERRORS}): the remainder is in the relay under its derived
id and, as far as the daemon knows, still in the outbox.  The next fetch
reconciles the two (the $ReadFetch$ replacement), and a second refile of the
same message is the same wire id, which the relay takes as one message
(\texttt{Nats-Msg-Id} within the stream's duplicate window; the local relay
compares ids).

What the bound says about this state.  The witness is the goal
$refileTimedOut = ztrue \land cancelled = zfalse \land mailOf~relay \setminus mailOf~held \neq
mailOf~relay$, a message in both stores right after the timeout.  With
$MAIL = 1 \upto 2$ and $fetchMax = 2$ the relay never holds more wires than one
fetch takes, so the next $ReadFetch$ always retrieves the surviving copy and
replaces the held original with it (the at-least-once rule working, not a
loss), and no state has the copy still in the relay once the fetch has run.  A
larger inbox than one fetch could leave the copy behind, and the state would
then persist across reads; the bound hides that, and the code's protection
there is that the refile is idempotent by id (the copy's id is derived from the
original's, so filing it again is the same message, never a second copy), and
that $Outbox.merge$ reconciles the two whenever the copy is finally fetched.

\begin{schema}{ReadCancelled}
  \Delta Store \\
  \Xi Client \\
  \Xi Ledger \\
  \Delta Flags \\
  rest? : \power WIRE
  \where
  phase = rdFetched \\
  rest? \subseteq held \\
  rest? \neq \emptyset \\
  relay' = relay \cup refile \limg rest? \rimg \\
  held' = held \\
  sent' = sent \setminus rest? \\
  phase' = rdIdle \\
  cancelled' = ztrue \\
  refileTimedOut' = refileTimedOut \\
  faulted' = faulted \\
  moreOut' = moreOut \\
  moreTruth' = moreTruth \\
  moreSlack' = moreSlack
\end{schema}

The read is cancelled after some of its refiles landed.  The connection's
closing cancels its running requests (\texttt{Connection.\_serve},
\texttt{task.cancel()}), so this is reachable.  It is a fault for the claim that
the remainder is reachable (what was not yet filed stays in the outbox only),
but not for the claim that no message is duplicated: the outbox releases each
message as it is filed (\texttt{Methods.\_refile}), so a message is in the
relay or in the outbox, never both, whenever the cancellation lands.  The superseded code released the filed messages once,
after the whole loop; that is the frozen variant
\texttt{mail-refile-cancel-dup.tex}, whose $ReadCancelled$ keeps $held' =
held$.

\begin{schema}{ClientShow}
  \Xi Store \\
  \Delta Client \\
  \Xi Ledger \\
  \Xi Flags \\
  w? : WIRE
  \where
  w? \in sent \\
  shown' = shown \cup \{ first~w? \} \\
  done' = done \\
  dupShown' = dupShown
\end{schema}

The client prints a message of the reply it holds.

\begin{schema}{Confirm}
  \Delta Store \\
  \Delta Client \\
  \Xi Ledger \\
  \Xi Flags \\
  w? : WIRE
  \where
  w? \in sent \\
  first~w? \in shown \\
  held' = held \setminus \{ w? \} \\
  sent' = sent \setminus \{ w? \} \\
  relay' = relay \\
  phase' = phase \\
  shown' = shown \\
  done' = done \cup \{ first~w? \} \\
  dupShown' = dupShown
\end{schema}

The client confirms a message, and the daemon marks it read and forgets it.
The guard $first~w? \in shown$ is the ordering of \texttt{Answer.deliver}:
print, flush, then confirm.

\begin{schema}{SessionEnd}
  \Delta Store \\
  \Xi Client \\
  \Delta Ledger \\
  \Xi Flags
  \where
  phase = rdIdle \\
  relay' = \emptyset \\
  held' = \emptyset \\
  sent' = \emptyset \\
  phase' = rdIdle \\
  issued' = issued \\
  drained' = drained \cup mailOf~relay \\
  lostBatch' = lostBatch \cup ( mailOf~sent \setminus shown ) \\
  lostRest' = lostRest \cup ( mailOf~( held \setminus sent ) \setminus shown )
\end{schema}

The session ends.  The drain moves what the relay holds into the owner's
mailbox, and the session's record, with its outbox, is discarded.

\begin{schema}{SessionEndDuringRead}
  \Delta Store \\
  \Xi Client \\
  \Delta Ledger \\
  \Delta Flags
  \where
  phase = rdFetched \\
  relay' = \emptyset \\
  held' = \emptyset \\
  sent' = \emptyset \\
  phase' = rdIdle \\
  issued' = issued \\
  drained' = drained \cup mailOf~relay \\
  lostBatch' = lostBatch \cup ( mailOf~sent \setminus shown ) \\
  lostRest' = lostRest \cup ( mailOf~( held \setminus sent ) \setminus shown ) \\
  cancelled' = cancelled \\
  refileTimedOut' = refileTimedOut \\
  faulted' = ztrue \\
  moreOut' = moreOut \\
  moreTruth' = moreTruth \\
  moreSlack' = moreSlack
\end{schema}

The session ends between a read's fetch and its reply.

\begin{schema}{DaemonCrash}
  \Delta Store \\
  \Xi Client \\
  \Delta Ledger \\
  \Delta Flags
  \where
  relay' = relay \\
  held' = \emptyset \\
  sent' = \emptyset \\
  phase' = rdIdle \\
  issued' = issued \\
  drained' = drained \\
  lostBatch' = lostBatch \cup ( mailOf~sent \setminus shown ) \\
  lostRest' = lostRest \cup ( mailOf~( held \setminus sent ) \setminus shown ) \\
  cancelled' = cancelled \\
  refileTimedOut' = refileTimedOut \\
  faulted' = ztrue \\
  moreOut' = moreOut \\
  moreTruth' = moreTruth \\
  moreSlack' = moreSlack
\end{schema}

The daemon dies.  The relay keeps what it holds; the outbox is gone.

\begin{schema}{Quiesce}
  \Xi MailState
\end{schema}

The stutter that stands for ``nothing more happens'': the bound on
$MAIL$ ends the history of arrivals, and without $Quiesce$ ProB would
report the end of the bounded history as a deadlock.

%%----------------------------------------------------------------------
\section{Model Checking}
%%----------------------------------------------------------------------

$MailInv$ is the state invariant; probcli reports it for the shipped
operations in the transcript at the end of this file.

These things the specification does \emph{not} promise, each reachable and
each stated as a reachability formula so that it is not an omission:

\begin{enumerate}
  \item \textbf{A duplicate delivery is possible.}  $EF\;(dupShown = ztrue)$
    holds, along $Arrive \to ReadFetch \to ReadReply \to ClientShow \to
    ReadFetch \to ReadReply$: a client that prints, dies before it confirms,
    and reads again is shown the message again, because the batch is still in
    the outbox and the next read offers it.  Delivery is at least once.  The same read also
    refiles a message that an earlier, unconfirmed reply already showed.
  \item \textbf{A batch is lost if the daemon dies or the session ends
    before the client has been shown it.}  $EF\;(lostBatch \neq \emptyset)$
    holds, along the trace $Arrive \to ReadFetch \to ReadReply \to
    SessionEnd$ with no $ClientShow$: the batch of an unconfirmed reply is
    held only in the outbox, and neither the drain nor a restart reads the
    outbox.  A client that dies between the reply and its print keeps the
    message unread \emph{only while the session and the daemon live}.  Two
    separate obligations close this.  (a) Session end routed through the
    daemon refiles the held batch before it drops the record: a hard
    requirement of the A2 change (DES-099), after which $SessionEnd$ no longer
    reaches $lostBatch \neq \emptyset$.  (b) The daemon-crash window is
    PL-1753: acknowledge the relay only after the client has shown the mail.
  \item \textbf{A refile that landed but timed out can deliver a message
    twice.}  $ReadReplyTimedOut$ leaves the filed copy in the relay and the
    original held; the next fetch reconciles them, so the two are never offered
    together ($held$ holds one copy of each message, an invariant).  But if the
    original is shown and confirmed first, the copy still in the relay is shown
    again; $EF\;(dupShown = ztrue \land faulted = ztrue)$ is reachable.
    This is the at-least-once duplicate, allowed.
  \item \textbf{After a timed-out refile the reply's \texttt{more} can
    overcount, by a stated bound.}  The remainder is still held and its landed
    copy is in the relay, and the code adds the held remainder to the relay's
    own count, so each such message is counted twice.  The invariant is not
    switched off; it is replaced by the exact statement the code satisfies:
    $moreTruth \leq moreOut \land moreOut \leq moreTruth + moreSlack$.  In plain
    words, \texttt{more} never undercounts the messages still unread, and it
    overcounts by at most $moreSlack$, the number of held messages whose filed
    copy is already in the relay (the timed-out-but-landed refiles outstanding
    at that reply).  Where $moreSlack = 0$ the count is exact.  The overcount is
    reachable ($EF\;(moreOut > moreTruth)$) and the bound is checked over the
    whole state space.  The next read reconciles it: the copy's id is derived
    from the original's, so the fetched copy is the same message, $ReadFetch$
    replaces the held original with it, and the next reply counts it once.
  \item \textbf{Order across reads is not guaranteed.}  A refile appends to
    the relay, behind mail that arrived during the read, so a later read can
    show a newer message before an older refiled one.  Within one reply the mail
    is ordered by timestamp, ties by id.  The model has no clock and does not
    state this; it is an ordering limit, not a loss.
  \item \textbf{The fetch-to-refile window is a crash window.}  A consuming
    fetch removes the mail from the relay before the remainder is filed
    again; $DaemonCrash$ between $ReadFetch$ and $ReadReply$ puts the
    remainder in $lostRest$ ($EF\;(lostRest \neq \emptyset \land faulted =
    ztrue)$).  The Relay API has no acknowledge-after-write, so this is the
    same obligation as (b), PL-1753, and no more than that is claimed here.
\end{enumerate}

Three frozen variants are the negative controls.
\texttt{mail-refile-cancel-dup.tex} has a $ReadCancelled$ that leaves the
filed messages in the outbox (the release after the loop): the goal $cancelled
= ztrue \land refileTimedOut = zfalse \land mailOf~relay \setminus mailOf~held
\neq mailOf~relay$ is reached right after the cancellation (a message in the
relay and in the outbox), where the shipped model does not reach it, and
$MailInv$ is violated at the next fetch.
\texttt{mail-refile-drop.tex} has a $ReadReply$ that files nothing back
(the code before the refile): $MailInv$ is violated, with the trace $Arrive
\to Arrive \to ReadFetch \to ReadReply \to SessionEnd$, the remainder lost
with no fault.  \texttt{mail-refile-confirm-early.tex} has a $Confirm$ with no
$first~w? \in shown$ guard (the confirmation sent before the print):
$MailInv$ is violated by a message marked read and never shown.

%%----------------------------------------------------------------------
%% Reproducing the checks.  Run from the repo root; one probcli at a time.
%%
%%   PROBCLI="${PROBCLI:-$HOME/Applications/ProB/probcli}"
%%   P="-p DEFAULT_SETSIZE 1 -p MAXINT 6 -p MAX_INITIALISATIONS 20 \
%%      -p MAX_OPERATIONS 500 -p TIME_OUT 600000"
%%
%% The invariant, exhaustively: 3,979 states, 25,062 transitions, under a second.  With MAIL = 1..3 a run of the earlier model, which still reached
%% duplicate copies, exceeded 29 GB and was stopped at 11 minutes: MAIL = 1..3
%% is NOT checked.
%%
%%   ( ulimit -v 8000000; "$PROBCLI" docs/specs/mail-refile.tex $P -model_check )
%%
%% Reachability of the three limits (each reports GOAL found, shortest trace):
%%
%%   "$PROBCLI" docs/specs/mail-refile.tex $P -model_check \
%%       -goal "lostBatch/={} & faulted=zfalse"
%%   "$PROBCLI" docs/specs/mail-refile.tex $P -model_check \
%%       -goal "dupShown=ztrue & faulted=zfalse"
%%   (a message in both stores, each event isolated by its own flag:)
%%   "$PROBCLI" docs/specs/mail-refile.tex $P -model_check \
%%       -goal "cancelled=ztrue & refileTimedOut=zfalse & mailOf(relay)-mailOf(held) /= mailOf(relay)"
%%       (UNREACHABLE here; the same goal on mail-refile-cancel-dup.tex is
%%        reached by Arrive, ReadFetch, ReadCancelled)
%%   "$PROBCLI" docs/specs/mail-refile.tex $P -model_check \
%%       -goal "refileTimedOut=ztrue & cancelled=zfalse & mailOf(relay)-mailOf(held) /= mailOf(relay)"
%%       (reached by Arrive, Arrive, ReadFetch, ReadReplyTimedOut)
%%   "$PROBCLI" docs/specs/mail-refile.tex $P -model_check \
%%       -goal "moreOut>0 & faulted=zfalse"
%%
%% Negative controls, each of which must report an invariant violation:
%%
%%   "$PROBCLI" docs/specs/mail-refile-drop.tex $P -model_check
%%   "$PROBCLI" docs/specs/mail-refile-confirm-early.tex $P -model_check
%%   "$PROBCLI" docs/specs/mail-refile-cancel-dup.tex $P -model_check
%%----------------------------------------------------------------------

\end{document}
